\documentclass[11pt,a4paper]{article}

\usepackage[T1]{fontenc}
\usepackage[utf8]{inputenc}
\usepackage{lmodern}
\usepackage[a4paper,margin=1in]{geometry}
\usepackage{amsmath}
\usepackage{booktabs}
\usepackage{graphicx}
\usepackage{float}
\usepackage{microtype}
\usepackage[hidelinks]{hyperref}

\title{Enhancing 6D Object Pose Estimation}
\author{TODO: Add author name and university details}
\date{\today}

\begin{document}

\maketitle

\begin{abstract}
This report investigates 6D object pose estimation for a LineMOD object using
an RGB-only model and a simple RGB--D feature-fusion model. The implemented
pipeline first uses YOLO for object localization and then predicts object
translation and rotation represented by a quaternion. The RGB--D model combines
features from an RGB ResNet-50 branch and a lightweight depth branch. The saved
experiments report object detection performance using mAP and pose performance
using the ADD metric. The experimental comparison is limited to the trained
checkpoints and LineMOD object 1 test split available in this repository.
\end{abstract}

\section{Introduction}

Six-degree-of-freedom (6D) object pose estimation seeks to estimate an
object's three-dimensional translation and orientation relative to a camera.
This information is useful for applications such as robotic manipulation,
augmented reality, and vision-based inspection. The task is more demanding
than two-dimensional detection because the output must describe the object's
position and orientation in three dimensions.

This project studies a modular pipeline for 6D pose estimation using the
LineMOD dataset. The implemented pipeline combines object detection with pose
regression. It also compares an RGB-only estimator with an RGB--D extension
that uses aligned depth information. The goal is to evaluate whether the
additional depth branch improves the observed pose accuracy for the experiment
stored in the repository.

The contributions represented by the current implementation are:
\begin{itemize}
    \item a LineMOD data-loading and preprocessing pipeline for RGB, depth, and
    pose annotations;
    \item YOLO-based object localization using LineMOD bounding-box labels;
    \item RGB pose regression with an ImageNet-initialized ResNet-50 backbone;
    \item a compact RGB--D feature-fusion model; and
    \item quantitative comparison using object detection metrics and the ADD
    pose metric.
\end{itemize}

\section{Related Work}

Object detection and 6D pose estimation are commonly treated as related but
separate stages. A detector provides an image region in which a pose estimator
can regress the object's translation and orientation. RGB--D methods extend
this setting by using depth to provide geometric information that is not
explicitly available in RGB intensity values.

The present project is inspired by this general RGB--D fusion approach, but it
implements a deliberately compact architecture suitable for a university
project rather than a full reproduction of a large pose-estimation system.

% TODO: Add verified scholarly references and a literature comparison. No
% citations are included here because they are not stored in the repository.

\section{Dataset and Data Preparation}

The project uses the preprocessed LineMOD dataset. The repository contains
RGB images, 16-bit depth images, object bounding boxes, ground-truth camera-to-
object rotations and translations, camera intrinsics, object masks, and PLY
models. The current saved experiment uses object 1.

For object 1, the repository's split files contain 186 training samples and
1050 test samples. The reported pose experiments focus on these object-1
splits.

The dataset loader reads the object-specific train and test split files and
selects the annotation corresponding to the configured object identifier. The
RGB and depth images are cropped using the same ground-truth object bounding
box and resized to $224\times224$ pixels. RGB inputs are converted to tensors
and normalized using ImageNet channel statistics. Depth is resized with nearest
neighbour interpolation and converted to metres for neural-network input.
Ground-truth translations and PLY model coordinates remain in millimetres,
which is also the unit used when reporting ADD.

The pose annotations provide a $3\times3$ rotation matrix and a three-
dimensional translation vector. The implemented geometry utility converts the
rotation matrix to a quaternion in $[w,x,y,z]$ order. The models normalize
predicted quaternions before returning them.

\section{Methodology}

\subsection{Object Detection}

The detection stage uses the pretrained YOLO11n checkpoint stored by the
repository and is fine-tuned on a YOLO-format version of the LineMOD object
bounding-box annotations. The dataset preparation script converts each
bounding box from $[x,y,w,h]$ image coordinates to the normalized YOLO
representation. Training is configured for 20 epochs with batch size 8 and an
input image size of 640 pixels. The saved YOLO configuration records random
seed 0. Detection performance is evaluated using mAP@50 and mAP@50:95; the
result values are reported in Section~\ref{sec:experimental-results}.

The repository also contains an end-to-end inference path in
\texttt{run\_pose\_pipeline.py}. In that path, YOLO first localizes the object
and the resulting crop is passed to either the RGB or RGB--D pose model.

\subsection{RGB Pose Estimation}

The RGB pose estimator in \texttt{src/pose\_model.py} uses a ResNet-50
initialized with ImageNet pretrained weights. Its classification layer is
replaced by an identity mapping, leaving the backbone feature vector as the
input to two linear regression heads. The translation head predicts
$(x,y,z)$, and the rotation head predicts a four-dimensional quaternion in
$[w,x,y,z]$ order. The predicted quaternion is normalized to unit length.

Training is configured for 20 epochs with batch size 4 using the Adam
optimizer and learning rate $10^{-4}$. Twenty percent of the original
training split is used as a validation subset, with random split seed 42.
Checkpoint selection is based on validation ADD, after which the selected
checkpoint is evaluated on the untouched test split.

\subsection{RGB--D Pose Estimation and Fusion}

The RGB--D model in \texttt{src/rgbd\_pose\_model.py} contains two input
branches. The RGB branch is a ResNet-50 initialized with ImageNet weights; its
final classification layer is removed and the resulting 2048-dimensional
feature vector is used for fusion. The depth branch is a lightweight CNN with
convolutional channel dimensions $1\rightarrow32\rightarrow64\rightarrow128$.
Each convolution uses stride 2 and is followed by a ReLU and batch-normalization
layer. Adaptive average pooling produces a 128-dimensional depth feature.

The 2048-dimensional RGB feature and 128-dimensional depth feature are
concatenated. The fusion module maps the resulting 2176-dimensional vector to
512 dimensions, applies ReLU and dropout with probability 0.3, and then maps
it to 256 dimensions with another ReLU. Separate linear heads predict the
three-dimensional translation and the four-dimensional quaternion rotation;
the quaternion is normalized to unit length.

Training is configured for 20 epochs with batch size 4 and learning rate
$10^{-4}$. The validation fraction is 0.2 and the random split seed is 42.
Checkpoint selection is based on validation ADD, followed by final evaluation
on the untouched test split.

This design is a compact RGB--D fusion implementation for the project and does
not claim to reproduce the complete architecture or training procedure of any
particular external RGB--D pose-estimation system.

\subsection{Loss Function and ADD Evaluation Metric}

The training loss combines translation mean squared error with a quaternion
rotation loss. Since $q$ and $-q$ represent the same rotation, the rotation
term uses the absolute quaternion dot product:
\begin{equation}
    \mathcal{L}_{\mathrm{rot}}
    = 1 - \left|q_{\mathrm{pred}}^{\mathsf{T}}q_{\mathrm{target}}\right|.
\end{equation}
The total loss is the weighted sum of the translation and rotation terms. The
training scripts instantiate the loss with its default weights, so both the
translation and rotation terms have weight 1.0:
\begin{equation}
    \mathcal{L}_{\mathrm{total}}
    = \mathcal{L}_{\mathrm{trans}} + \mathcal{L}_{\mathrm{rot}}.
\end{equation}

For evaluation, ADD transforms the object's 3D model points using the predicted
and ground-truth poses and computes the mean Euclidean distance between the two
sets of transformed points:
\begin{equation}
    \mathrm{ADD}
    = \frac{1}{N}\sum_{i=1}^{N}
    \left\|\left(x_iR_{p}^{\mathsf{T}}+t_{p}\right)
    -\left(x_iR_{g}^{\mathsf{T}}+t_{g}\right)\right\|_2.
\end{equation}
The repository stores the model points and translations in millimetres, so the
reported ADD values are also in millimetres.

The project uses the preprocessed LineMOD dataset. The reported pose
experiments focus on object 1, which contains 186 training samples and 1050
test samples. The dataset provides RGB images, depth images, masks, object
bounding boxes, camera information, 6D pose annotations, and PLY object models
for the available objects.

For pose training and quantitative pose evaluation, the RGB and RGB--D dataset
loaders use the ground-truth object bounding boxes to create the input crops.
The pose-network inputs are resized to $224\times224$ pixels. RGB inputs use
ImageNet normalization, while depth is resized using the repository's
nearest-neighbour preprocessing and converted to metres for network input.
Translation targets, object-model coordinates, pose errors, and ADD results
remain expressed in millimetres.

The saved experimental comparison concerns LineMOD object 1. Pose results are
reported on the untouched test split, while the saved YOLO detection result is
reported for the detector validation output. Checkpoints are loaded from the
repository's \texttt{checkpoints} directory, and final metrics are read from
the machine-readable files in \texttt{outputs} and the saved YOLO run results.

The reported RGB versus RGB--D ADD values were computed using ground-truth
object crops from the dataset loaders. They were not computed from YOLO-
detected crops. Separately, \texttt{run\_pose\_pipeline.py} demonstrates an
end-to-end inference path in which YOLO localization is followed by pose
estimation; detector-crop ADD was not used for the reported quantitative
comparison.

The final results section uses only values that are already stored in the
repository. No additional training or evaluation is performed by this report
source file.

\section{Experimental Results}\label{sec:experimental-results}

\subsection{Object Detection Results}

The YOLO detector was evaluated on the LineMOD object 1 validation split. The best recorded row in the saved YOLO results contains an mAP@50 of $0.99500$ and an mAP@50:95 of $0.89686$. These values are reported directly from the repository's saved detection results.

\begin{table}[htbp]
    \centering
    \caption{YOLO object detection results for LineMOD object 1.}
    \label{tab:yolo-results}
    \begin{tabular}{lc}
        \hline
        Metric & Value \\
        \hline
        mAP@50 & 0.99500 \\
        mAP@50:95 & 0.89686 \\
        \hline
    \end{tabular}
\end{table}

\subsection{Pose Estimation Comparison}

The pose estimators were evaluated on the LineMOD object 1 test split using the ADD metric in millimetres. The RGB-only model obtained an ADD of $589.5225\,\mathrm{mm}$, while the RGB--D model obtained an ADD of $138.2051\,\mathrm{mm}$. Because lower ADD indicates a smaller average distance between the transformed model points, the RGB--D result was lower for this experiment.

The percentage change was calculated from the saved ADD values as
\[
\frac{589.5225 - 138.2051}{589.5225} \times 100 = 76.5564\%.
\]
Thus, the RGB--D experiment produced a verified $76.5564\%$ reduction in ADD relative to the RGB-only experiment on this test split. This observation applies to the reported object 1 experiment and does not establish that RGB--D input is universally superior in other datasets, objects, or experimental settings.

\begin{table}[htbp]
    \centering
    \caption{Comparison of RGB-only and RGB--D pose estimation on the LineMOD object 1 test split.}
    \label{tab:pose-results}
    \begin{tabular}{lcc}
        \hline
        Metric & RGB & RGB--D \\
        \hline
        ADD (mm) & 589.5225 & 138.2051 \\
        \hline
    \end{tabular}
\end{table}

\begin{figure}[htbp]
    \centering
    \includegraphics[width=0.78\textwidth]{outputs/rgb_vs_rgbd_add.png}
    \caption{ADD comparison between the RGB-only and RGB--D pose estimators.}
    \label{fig:rgb-rgbd-add}
\end{figure}

\begin{figure}[htbp]
    \centering
    \includegraphics[width=\textwidth]{outputs/pose_prediction_examples.png}
    \caption{Qualitative pose prediction examples for LineMOD object 1 test samples. The figure reports the ground-truth translation, RGB and RGB--D predicted translations, translation errors, and quaternion angular errors.}
    \label{fig:pose-prediction-examples}
\end{figure}

Overall, the saved results show substantially lower ADD for the RGB--D model than for the RGB-only model in this experiment. The qualitative examples provide sample-level context for the reported predictions and errors, while the numerical ADD values provide the aggregate comparison. The conclusions are limited to the evaluated LineMOD object 1 test split and the trained checkpoints represented by the saved repository outputs.


\section{Discussion}

The saved experiment shows a lower ADD value for the RGB--D model than for the
RGB-only model on the evaluated LineMOD object 1 test split. The detection
results also show high recorded YOLO mAP values for the saved validation
output. Together, these results indicate that the implemented detector and
pose-estimation pipeline produced useful performance for this specific
experiment.

The comparison should be interpreted within its stated scope. It concerns one
LineMOD object, the available trained checkpoints, and the recorded evaluation
splits. It does not establish that RGB--D input is universally better than RGB
input, nor does it replace a broader comparison across objects, runs, datasets,
or controlled ablation settings.

% TODO: Add an error analysis based on additional verified outputs, such as
% failure cases, depth-quality analysis, or per-sample statistics.

\section{Conclusion and Future Work}

This project implements a complete experimental pipeline from LineMOD data
preparation and YOLO localization to RGB and RGB--D 6D pose estimation. The
models predict translation and quaternion rotation, and the evaluation code
reports ADD using the object's 3D model. The saved object 1 experiment provides
both quantitative and qualitative evidence for comparing RGB-only and RGB--D
pose prediction.

Future work should include broader evaluation across the available LineMOD
objects, repeated runs with controlled random seeds, and additional analysis of
failure cases. Other possible extensions include improved depth handling,
stronger augmentation, and a more systematic ablation of the RGB--D fusion
components.

% TODO: Replace future-work suggestions with the final experimentally justified
% scope of the project.

\section*{References}
\addcontentsline{toc}{section}{References}

% TODO: Add the verified bibliographic entries required by the final report.
% No citations are included until their sources and bibliographic details have
% been confirmed.

\end{document}
